Seien $\alpha_1$, $\alpha_2$ und $\alpha_3$ die drei verschiedenen
Nullstellen des irreduziblen Polynoms
\[
m(X)
=
X^3 + 2X^2 + 9X + 1.
\]
Berechnen Sie 
\[
S_k
=
\alpha_1^k
+
\alpha_2^k
+
\alpha_3^k
\]
für $k=1,2,3$.

\begin{loesung}
Die Polynome $X_1^k+X_2^k+X_3^k$ sind symmetrisch und müssen sich
daher durch die elementarsymmetrischen Polynome ausdrücken lassen.
Die elementarsysmmetrischen Polynome von drei Variablen sind in
\eqref{buch:komplex:fundamental:bsp:elementar3}
berechnet worden.
Zur Vereinfachung der Notation lassen wir die Argument $X_k$ weg.
Setzt man die Nullstellen in die elemetarsymmetrischen Funktionen
ein, ergeben sich die Koeffizienten des Polynoms $m(X)$:
\[
\sigma_1 = -2,\qquad
\sigma_2 = 9\qquad\text{und}\qquad
\sigma_3 = -1.
\]

Für $k=1$ ist $S_1 = \sigma_1(\alpha_1,\alpha_2,\alpha_3)$,
dies ist der negative Koeffizient von $X^2$ des Polynoms $m(X)$,
also $S_1 = -2$.

Für $k=2$ müssen wir erst einen Ausdruck $S_2$ finden:
\begin{align*}
\sigma_1^2
&=
\sum_{i=1}^3X_i^2 + \sum_{i\ne j}^3 X_iX_j
=
\sum_{i=1}^3X_i^2 + 2\sigma_2
\intertext{oder aufgelöst nach der gesuchten Summe}
\sum_{i=1}^3X_i^2
&=
\sigma_1^2
- 2\sigma_2.
\end{align*}
Durch Einsetzen der Nullstellen finden wir
\begin{align*}
S_2 &= 
\sigma_1(\alpha_1,\alpha_2,\alpha_3)^2
- 2\sigma_2(\alpha_1,\alpha_2,\alpha_3)
=
(-2)^2-2\cdot 9
=-14.
\end{align*}

Für $k=3$ können wir das Resultat~\eqref{buch:komplex:symmetrisch:eqn:X33}
\begin{align*}
\sum_{i=1}^3 X_i^3
&=
\sigma_1^3-3\sigma_1\sigma_2+3\sigma_3
\intertext{verwenden.
Einsetzen der Nullstellen ergibt}
S_3
=
\sum_{i=1}^3\alpha_i^3
&=
(-2)^3 - 3\cdot (-2)\cdot 9 + 3\cdot(-1)
=
-8+54-3=43.
\qedhere
\end{align*}
\end{loesung}
